\documentclass[11pt,a4paper]{article}
\usepackage[top=2.8cm,bottom=2.8cm,left=3cm,right=3cm]{geometry}
\usepackage[utf8]{inputenc}
\usepackage[T1]{fontenc}
\usepackage{lmodern}
\usepackage{microtype}
\usepackage{setspace}
\usepackage[table,dvipsnames]{xcolor}
\usepackage{tcolorbox}
\tcbuselibrary{skins,breakable,hooks}
\usepackage{booktabs}
\usepackage{longtable}
\usepackage{array}
\usepackage{graphicx}
\usepackage{hyperref}
\usepackage{parskip}
\usepackage{titlesec}
\usepackage{enumitem}
\usepackage{fancyhdr}
\usepackage{amsmath}
\usepackage{amssymb}

\definecolor{primary}{HTML}{1A1A2E}
\definecolor{accent}{HTML}{C0392B}
\definecolor{highlight}{HTML}{1A5276}
\definecolor{linkblue}{HTML}{154360}
\definecolor{boxbg}{HTML}{F8F9FA}
\definecolor{termbg}{HTML}{EBF5FB}
\definecolor{industrybg}{HTML}{EAFAF1}
\definecolor{trapbg}{HTML}{FDEDEC}
\definecolor{thinkbg}{HTML}{FEF9E7}
\definecolor{curiobg}{HTML}{F4ECF7}
\definecolor{insightbg}{HTML}{EBF5FB}

\hypersetup{colorlinks=true,linkcolor=linkblue,urlcolor=linkblue,citecolor=linkblue}
\setstretch{1.15}

\newtcolorbox{termbox}[1]{title={\textbf{Term: #1}},colback=termbg,colframe=highlight,coltitle=white,fonttitle=\bfseries\small,breakable,skin=enhanced}
\newtcolorbox{industrybox}[1]{title={\textbf{Industry: #1}},colback=industrybg,colframe=green!50!black,coltitle=white,fonttitle=\bfseries\small,breakable,skin=enhanced}
\newtcolorbox{trapbox}[1]{title={\textbf{Trap: #1}},colback=trapbg,colframe=accent,coltitle=white,fonttitle=\bfseries\small,breakable,skin=enhanced}
\newtcolorbox{thinkbox}{title={\textbf{Think About It}},colback=thinkbg,colframe=orange!70!black,coltitle=white,fonttitle=\bfseries\small,breakable,skin=enhanced}
\newtcolorbox{curiobox}[1]{title={\textbf{Q: #1}},colback=curiobg,colframe=purple!60!black,coltitle=white,fonttitle=\bfseries\small,breakable,skin=enhanced}
\newtcolorbox{insightbox}{title={\textbf{Key Insight}},colback=insightbg,colframe=highlight,coltitle=white,fonttitle=\bfseries\small,breakable,skin=enhanced}
\newtcolorbox{quizbox}[1]{title={\textbf{Question #1}},colback=boxbg,colframe=primary,coltitle=white,fonttitle=\bfseries\small,breakable,skin=enhanced}
\newtcolorbox{answerbox}[1]{title={\textbf{Answer #1}},colback=industrybg,colframe=green!50!black,coltitle=white,fonttitle=\bfseries\small,breakable,skin=enhanced}

\pagestyle{fancy}
\fancyhf{}
\fancyhead[L]{\small\textcolor{primary}{L1-14: AI/ML Infrastructure Basics}}
\fancyhead[R]{\small\textcolor{primary}{System Design Curriculum}}
\fancyfoot[C]{\thepage}
\renewcommand{\headrulewidth}{0.4pt}

\titleformat{\section}{\large\bfseries\color{primary}}{}{0em}{}[\titlerule]
\titleformat{\subsection}{\normalsize\bfseries\color{highlight}}{}{0em}{}

\begin{document}

%--- TITLE PAGE ---
\begin{titlepage}
\centering
\vspace*{2cm}
{\color{primary}\rule{\linewidth}{1.5pt}}\\[0.4cm]
{\huge\bfseries\color{primary} L1-14: AI/ML Infrastructure Basics}\\[0.3cm]
{\Large\color{highlight} Feature Stores, Vector Databases, Embeddings, Model Serving, and Inference Architecture}\\[0.2cm]
{\color{primary}\rule{\linewidth}{1.5pt}}\\[1cm]
{\large System Design Curriculum --- Layer 1: Foundations}\\[0.4cm]
{\normalsize Edition 2025}\\[0.3cm]
{\small Data sources: Netflix Tech Blog, Spotify Engineering Blog, AWS re:Invent 2023--2024,\\
ANN-Benchmarks 2025, GitGuardian 2025, CNCF Annual Survey 2024, QCon 2023 Spotify Hendrix talk}\\[2cm]

\begin{insightbox}
In 2025, senior engineers at Apple, Google, Netflix, and Spotify are expected to understand the infrastructure that \textit{serves} models in production---not the math that trains them. This lesson is about the plumbing: how features are computed and served with consistency, how embeddings enable similarity search at billion-vector scale, and what inference latency means for system design tradeoffs. A candidate who cannot speak to model serving is missing a critical architectural layer in every modern recommendation, fraud detection, or personalization system.

The signal interviewers look for: do you understand \textit{why} a feature store exists (training-serving skew), what a vector database actually solves (ANN search), and what the latency constraints are (recommendation $<$100ms, fraud detection $<$50ms)? This lesson answers all three.
\end{insightbox}

\vfill
{\small \textit{System Design Interview Curriculum} --- Lesson L1-14 (Final Layer 1 Lesson)}
\end{titlepage}

\tableofcontents
\addtocontents{toc}{\protect\setcounter{tocdepth}{2}}
\newpage

%--- SECTION 1: WHY THIS EXISTS ---
\section{Why This Exists}

Modern consumer software is built on machine learning predictions. Netflix's recommendation system influences 80\% of the content watched on the platform. Spotify's Discover Weekly playlist, generated weekly for over 600 million users, uses collaborative filtering over hundreds of millions of listening histories. Uber's fraud detection system scores 20 million trips per day, producing a risk signal in under 50 milliseconds. None of these systems work without infrastructure: pipelines to compute features, stores to serve them consistently, model servers to run inference, and latency budgets tight enough not to degrade the user experience.

The engineering challenge is that ML infrastructure lives at the intersection of two conflicting paradigms. Batch systems (Spark, data warehouses) optimize for throughput over correctness, processing terabytes of data with eventual consistency. Real-time systems (Redis, in-memory caches) optimize for latency over throughput, serving requests in single-digit milliseconds. A recommendation system needs both: it trains on batch-processed historical data but serves predictions in real time. Bridging this gap---consistently, at scale, without the model seeing different data than it was trained on---is what ML infrastructure is designed to do.

The reason interviewers probe this in senior rounds is that ML systems have failure modes that traditional distributed systems do not. A database bug causes visible errors. A training-serving skew bug causes the model to silently serve wrong predictions---predictions that \textit{look right} but gradually degrade user experience. Detecting and preventing this class of silent correctness bug requires understanding the full ML pipeline from data collection through feature engineering through serving.

%--- SECTION 2: EMBEDDINGS AND VECTOR REPRESENTATIONS ---
\section{Embeddings and Vector Representations}

\subsection{What an Embedding Is}

An embedding is a dense, fixed-length vector of floating-point numbers that encodes semantic meaning. A document, a product, a user's listening history, or an image is mapped by a neural network to a point in a high-dimensional space---typically 128 to 3072 dimensions depending on the model. The key property learned during training is that semantically similar items land \textit{close together} in this space. Two songs from the same genre have embeddings whose cosine similarity is high; a jazz track and a death metal track are far apart. This geometry is what enables the core ML application: ``find me the K items most similar to this query item'' becomes a nearest-neighbor search problem in embedding space.

The practical implication for system design is that raw embeddings are vectors of floats: an OpenAI \texttt{text-embedding-3-small} embedding is 1536 dimensions $\times$ 4 bytes (float32) = 6,144 bytes per item. Storing 100 million product embeddings requires 600GB of raw float data. Querying that data with a brute-force nearest-neighbor search (compute distance to every vector) would take hundreds of milliseconds per query at that scale---far too slow for real-time serving. This is why approximate nearest-neighbor (ANN) algorithms and dedicated vector databases exist.

\begin{industrybox}{OpenAI Embedding Models 2024--2025}
OpenAI's \texttt{text-embedding-3-small} (1536 dimensions) and \texttt{text-embedding-3-large} (3072 dimensions) are the most widely used text embedding APIs as of 2024. The models support dimensionality reduction via the API---\texttt{text-embedding-3-small} can be truncated to 256 dimensions with minimal quality loss, enabling a 6x storage reduction for applications where 256 dimensions is sufficient. Embedding API usage grew substantially in 2023--2024 driven by RAG (Retrieval-Augmented Generation) pipelines: as LLMs became mainstream, the pattern of embedding a knowledge base and retrieving relevant chunks at inference time became the dominant architecture for grounding LLM responses in private data.
\end{industrybox}

\subsection{Approximate Nearest Neighbor Search}

Exact nearest-neighbor search in high-dimensional spaces is inherently expensive: the so-called ``curse of dimensionality'' means that as dimensions increase, all points become approximately equidistant and efficient spatial indexes (like k-d trees, which work well in 2--10 dimensions) degrade to linear scan. ANN algorithms trade a small loss in recall (some true nearest neighbors are missed) for dramatic speedups.

The dominant ANN algorithm in production systems is HNSW (Hierarchical Navigable Small World graphs), developed in 2016. HNSW builds a multi-layer graph where each vector is a node; the top layers have sparse long-range connections for fast global traversal, and the bottom layers have dense local connections for precise search. At query time, the algorithm navigates the graph greedily, reaching approximate nearest neighbors in O(log n) hops rather than O(n) linear scan. HNSW achieves 99\% recall at $<$10ms query latency for 100 million vectors with the right index parameters. The tradeoff is memory: HNSW stores the full index in RAM, which at 100M vectors with 1536 dimensions can reach hundreds of GB.

%--- SECTION 3: VECTOR DATABASES ---
\section{Vector Databases}

\subsection{Why a Dedicated Database}

Traditional relational and NoSQL databases store data in formats optimized for exact-match queries (B-tree indexes), range scans (sorted keys), or document retrieval (inverted indexes). None of these structures support efficient ANN search. PostgreSQL with the \texttt{pgvector} extension can perform vector search, but it is limited to exact nearest-neighbor (using sequential scan) or approximate search via HNSW indexes built in the extension. At under 1 million vectors, \texttt{pgvector} with HNSW is competitive with dedicated vector databases and has the advantage of eliminating an operational dependency. Above 10 million vectors, dedicated vector databases significantly outperform it in both query throughput and memory efficiency.

\begin{industrybox}{Vector Database Market 2024--2025}
The global vector database market reached \$3.2 billion in 2025, growing at over 24\% annually (market analyst data, 2025). The leading options by deployment pattern: Pinecone (managed, purpose-built, closed-source) raised over \$130 million and serves enterprise RAG and search workloads. Qdrant (open-source, written in Rust) benchmarks at 1840 QPS on 1M-vector workloads at ANN-Benchmarks 2025, the highest of major vector databases. Weaviate (open-source) exceeds 1 million Docker pulls per month and supports hybrid search combining vector similarity with keyword filters. \texttt{pgvector} adoption grew 3x in 2024 as teams embedded vector search in existing PostgreSQL infrastructure to avoid new operational dependencies.
\end{industrybox}

\subsection{When to Use What}

The decision between vector databases comes down to scale, team operational capacity, and query complexity. For teams with under 10 million vectors and an existing PostgreSQL deployment, \texttt{pgvector} eliminates infrastructure complexity. Above 10 million vectors, or when query latency requirements drop below 10ms at high QPS, dedicated vector databases provide better performance. Managed services (Pinecone, Weaviate Cloud) are appropriate when the team cannot operationalize vector infrastructure; self-hosted Qdrant or Weaviate are appropriate for teams with the operational capacity and cost sensitivity. A common pattern is starting with \texttt{pgvector} and migrating to a dedicated vector database when query latency becomes a bottleneck---the data migration is straightforward since vectors are just float arrays.

\begin{trapbox}{``A vector database replaces my search engine''}
Vector search and keyword search solve different problems. Vector search finds semantically similar items even without exact keyword matches (``comfortable running shoe'' finds results containing ``cushioned athletic footwear'' even without word overlap). Keyword search (Elasticsearch, OpenSearch) finds exact or fuzzy matches with strict relevance control. The modern best practice is \textit{hybrid search}: combine vector similarity with keyword filtering in a single query. Weaviate and Elasticsearch both support this natively. A product search system that uses only vector search will miss results when users search for specific product codes or exact brand names. A system using only keyword search will miss semantically related results. Use both.
\end{trapbox}

%--- SECTION 4: FEATURE STORES ---
\section{Feature Stores}

\subsection{The Training-Serving Skew Problem}

Feature stores exist because of a specific, high-impact failure mode called \textit{training-serving skew}: the phenomenon where a feature is computed differently during model training than during model serving, causing the model to receive systematically different input in production than the input it was trained on. The predictions look numerically valid but are subtly wrong. The model silently degrades.

A concrete example: suppose you train a fraud detection model using a feature ``number of transactions in the past 24 hours by this user.'' During training, you compute this from a historical database with a Spark job that runs a complete window aggregation. In production, you compute it from a Redis counter that resets at midnight UTC, not on a rolling 24-hour window. The feature values are not the same. The model was trained on rolling-window counts and is served midnight-reset counts. For users who transact across midnight, the feature values diverge systematically---and the model's predictions for those users are wrong.

Feature stores solve this by centralizing the feature computation logic: the \textit{same code} that computes a feature for training also computes it for serving. The feature definition (``rolling 24-hour transaction count, computed as the sum of events in a sliding window anchored to the current timestamp'') lives in one place. Training pipelines and serving pipelines both call the same logic.

\begin{termbox}{Feature Store Architecture}
A feature store has three components: (1) the \textbf{feature registry}---a catalog of feature definitions with metadata (owner, data type, computation logic, freshness requirement); (2) the \textbf{offline store}---a data warehouse or data lake (S3 + Parquet, BigQuery, Snowflake) that stores historical feature values for training dataset generation; (3) the \textbf{online store}---a low-latency key-value store (Redis, DynamoDB, Bigtable) that serves the most recent feature values at inference time with $<$10ms latency. The write path computes features in a batch or streaming pipeline and writes to both stores; the training path reads from the offline store; the serving path reads from the online store.
\end{termbox}

\begin{industrybox}{Uber Michelangelo Feature Store 2024}
Uber's Michelangelo ML platform, described in their engineering blog as managing over 1,000 production features across fraud detection, ETA prediction, and driver matching, pioneered the feature store concept at scale. Their architecture separates feature computation (Flink and Spark pipelines), offline storage (Hive/HDFS for training), and online storage (Cassandra for serving). The key design principle: features are versioned and immutable---a feature definition change creates a new version, never mutates the existing one, preserving reproducibility of past model training runs. Netflix's Genie platform and Spotify's Hendrix (their ML platform, described at QCon New York 2023) follow similar architectures with centralized feature registries and dual offline/online stores.
\end{industrybox}

%--- SECTION 5: MODEL SERVING ---
\section{Model Serving Architecture}

\subsection{Online Inference vs Batch Inference}

ML predictions are served via two distinct patterns. Online inference produces a prediction on demand, synchronously, within the latency budget of the calling system---typically 10--200ms. A recommendation API returning the top-10 items for a user's homepage uses online inference. Batch inference runs predictions over a large dataset in advance, storing results for later lookup. Pre-computing recommendations for all users nightly and storing them in Redis for instant serving is batch inference. Batch inference amortizes compute costs but produces stale predictions; online inference is fresh but computationally expensive per-request.

The choice between them is driven by staleness tolerance and query volume. A homepage recommendation can tolerate 15-minute staleness; a fraud score cannot. Ad ranking requires fresh predictions because the ad auction changes by the millisecond. Most production systems use a combination: a ``precomputed'' tier (batch inference results stored in a fast cache) for the common case, and a ``real-time'' tier (online inference) that overrides the precomputed result when recency is critical or the user has recent strong signals (just watched something, just made a purchase).

\subsection{Model Servers and Latency}

Model serving infrastructure has standardized around a few key systems. TorchServe (PyTorch's native serving framework), TensorFlow Serving, and NVIDIA Triton Inference Server handle classical ML and deep learning models. Triton supports dynamic batching---accumulating multiple inference requests into a single GPU batch call, dramatically improving throughput at the cost of added latency per request. For large language models, vLLM (open-source, 2023) became the dominant inference engine by implementing PagedAttention, a memory management technique that enables efficient KV cache sharing across concurrent requests, increasing GPU throughput by 2--4x compared to naive implementations.

\begin{industrybox}{Inference Latency Budgets in Production 2024}
Industry SLA data from published engineering blogs: Netflix homepage recommendation: $<$100ms p99 end-to-end (Netflix Tech Blog, 2023). Uber fraud detection: $<$50ms p99 per trip score (Uber Engineering, 2023). Spotify's Discover Weekly pre-computation: batch job running 4 hours weekly, not latency-sensitive (Spotify Engineering Blog, 2022). LLM inference: time-to-first-token (TTFT) is the primary UX metric---GPT-4-class models on A100 GPUs achieve 1--3 seconds TTFT; smaller 7B-parameter models serve first token in 200--500ms. LLM throughput (tokens/second after first token) ranges from 20--100 tokens/s depending on model size, hardware, and batch size. GPU cold start for a 70B-parameter model: 30--120 seconds to load weights from storage into GPU VRAM.
\end{industrybox}

\subsection{GPU vs CPU Inference}

The choice between GPU and CPU inference involves cost, latency, and model size tradeoffs. Large deep learning models (transformer-based recommendation models, LLMs, image models) require GPU inference for acceptable latency---a 7B-parameter LLM on CPU takes seconds per token; on a single A100 GPU it takes milliseconds. Smaller traditional ML models (gradient boosted trees for fraud detection, logistic regression for click-through rate) run efficiently on CPU, which costs 10--50x less per inference than GPU. A well-architected ML serving infrastructure uses GPU for deep learning models and CPU for classical ML, routing requests appropriately.

%--- SECTION 6: THE ML PIPELINE ---
\section{The Full ML Pipeline}

Understanding the full lifecycle of an ML model is essential for system design interviews, because each stage has infrastructure implications and each transition between stages is a potential source of skew or failure.

The pipeline begins with \textbf{data collection}: raw events (clicks, purchases, views, fraud signals) are logged to an event bus (Kafka) and land in a data lake (S3) via a streaming consumer. Upstream data quality failures here propagate silently through the entire pipeline---garbage in, garbage out is particularly insidious in ML because the model degrades gradually rather than crashing immediately.

\textbf{Feature engineering} transforms raw events into model-ready features. This is where training-serving skew is introduced if not managed carefully through a feature store. Features are stored in the offline store for training and the online store for serving. \textbf{Model training} runs on the offline store's historical feature data, typically using a distributed training framework (Spark MLlib for classical ML, Ray or Horovod for distributed deep learning). The trained model is serialized and stored in a \textbf{model registry} (MLflow, SageMaker Model Registry) with version metadata, evaluation metrics, and the training data snapshot used.

\textbf{Model deployment} follows a controlled rollout: shadow mode (new model runs in parallel but predictions are not served---only logged for comparison), then canary (5--10\% of traffic), then full rollout. A/B testing frameworks assign users deterministically to model versions based on hashing, ensuring consistent assignment per user across requests. Finally, \textbf{model monitoring} tracks prediction distributions, feature distributions, and downstream business metrics to detect \textit{model drift}---the gradual degradation of prediction quality as the real-world data distribution shifts away from the training distribution.

\begin{thinkbox}
Why does model drift happen even when the code doesn't change? A fraud detection model trained on 2023 transaction patterns will gradually lose accuracy as fraud patterns evolve in 2024. A recommendation model trained in summer will have different genre preferences than what users want in winter. The model itself is frozen; the world it is predicting changes. Monitoring the model's output distribution (are fraud score distributions shifting?) and periodically retraining on fresh data is the operational response. This is distinct from concept drift (the underlying relationship changes) vs covariate shift (the input distribution changes while the relationship is stable)---though both require retraining.
\end{thinkbox}

%--- SECTION 7: DECISION FRAMEWORK ---
\section{Decision Framework: When to Use What}

\begin{center}
\begin{tabular}{p{2.8cm}p{2.8cm}p{2.8cm}p{3.4cm}}
\toprule
\textbf{Dimension} & \textbf{pgvector} & \textbf{Dedicated Vector DB} & \textbf{Keyword Search (ES)} \\
\midrule
Max practical scale & $\sim$10M vectors & 1B+ vectors & Unlimited (text) \\
Query latency & $<$20ms ($<$5M) & $<$10ms at 100M & $<$50ms \\
Operational overhead & Low (existing PG) & Medium--High & Medium \\
Hybrid search & Limited & Native (Weaviate) & Native \\
Cost & Low & Medium--High & Medium \\
Best for & Startup, $<$10M vecs & Scale-up, RAG & Exact/fuzzy text \\
\bottomrule
\end{tabular}
\end{center}

\vspace{0.5em}

\textbf{Batch vs Online Inference decision:} Use batch if predictions are needed for all users on a schedule and staleness of 15min to 24h is acceptable (homepage recommendations, weekly playlists). Use online inference when predictions must reflect events in the last seconds or minutes (fraud detection, search ranking, real-time ad auction). Hybrid: precompute for the common case, online inference for users with recent strong signals.

\textbf{Feature store adoption trigger:} When the same feature is computed in more than one place (training pipeline AND serving code), adopt a feature store. The cost of skew---silent model degradation---is always higher than the cost of the feature store.

%--- SECTION 8: CRITICAL NUMBERS ---
\section{Critical Numbers Reference}

\begin{center}
\begin{tabular}{p{5cm}p{3cm}p{5cm}}
\toprule
\textbf{Metric} & \textbf{Value} & \textbf{Source / Year} \\
\midrule
OpenAI text-embedding-3-small dims & 1536 & OpenAI API docs, 2024 \\
OpenAI text-embedding-3-large dims & 3072 & OpenAI API docs, 2024 \\
Storage per 1536-dim float32 vector & 6 KB & Calculated \\
100M vectors storage (float32) & 600 GB & Calculated \\
HNSW recall at 100M vectors & $\sim$99\% & ANN-Benchmarks 2025 \\
ANN query latency (HNSW, 100M) & $<$10ms & ANN-Benchmarks 2025 \\
Qdrant QPS benchmark (1M vectors) & 1840 QPS & ANN-Benchmarks 2025 \\
pgvector practical ceiling & $\sim$10M vectors & Community data, 2024 \\
Online feature retrieval (Redis) & $<$5ms p99 & Industry standard \\
Recommendation inference SLA & $<$100ms p99 & Netflix Tech Blog, 2023 \\
Fraud detection SLA & $<$50ms p99 & Uber Engineering, 2023 \\
LLM first-token latency (7B on GPU) & 200--500ms & vLLM benchmarks, 2024 \\
GPU cold start (70B model) & 30--120 seconds & OpenAI/Anthropic ops data \\
LLM throughput (tokens/s) & 20--100 t/s & Hardware-dependent \\
Vector DB market size & \$3.2B (2025) & Market analysts, 2025 \\
\bottomrule
\end{tabular}
\end{center}

%--- SECTION 9: SYSTEM DESIGN APPLICATION ---
\section{System Design Application}

\subsection{Design a Recommendation System (Netflix / Spotify)}

A recommendation system requires embedding-based similarity search, feature stores, and model serving all working together. At the data layer, user interaction events (clicks, plays, purchases) stream through Kafka into the feature pipeline. A batch Spark job computes long-horizon features (30-day genre affinity, historical rating distribution) and writes to the offline store (S3/Parquet). A Flink streaming job computes short-horizon features (last session, last 15 minutes of activity) and writes to the online store (Redis).

At serving time, the recommendation API request triggers a feature retrieval call to the online store ($<$5ms), then an embedding lookup for the user's current embedding ($<$10ms from the vector database), then an ANN search across the item embedding index to find the top-1000 candidates, then a ranking model run (online inference) to score and rerank the 1000 candidates down to 10 displayed items. The end-to-end latency target is under 100ms p99. The A/B testing framework assigns users to model variants based on a hash of their user ID, ensuring consistent assignment per session.

\begin{industrybox}{Netflix Recommendation Infrastructure 2024}
Netflix's recommendation system influences 80\% of all content watched on the platform. Their architecture (Netflix Tech Blog, 2023--2024) uses a multi-stage retrieval and ranking pipeline: candidate generation (ANN search from user embedding), filtering (availability, already-watched), and ranking (ensemble of collaborative filtering, deep neural networks, and contextual models). Netflix is consolidating multiple specialized models into a unified multi-task ``Hydra'' architecture that handles homepage ranking, search ordering, and notification personalization simultaneously, reducing the number of models requiring separate infrastructure. Spotify's Hendrix platform (described at QCon New York 2023) standardized on TFX and Kubeflow for ML pipelines, serving over half of Spotify's internal ML practitioners from a centralized infrastructure.
\end{industrybox}

\subsection{Design a Fraud Detection System}

Fraud detection requires the tightest latency budget of any common ML application: under 50ms end-to-end, because it sits in the critical path of a payment transaction. The architecture uses a precomputed feature pipeline: a Flink streaming job maintains a Redis hash of per-user risk signals updated in real time (recent transaction velocity, geolocation change in the last 5 minutes, device fingerprint seen before). When a transaction arrives, the fraud scoring service reads all features from Redis ($<$5ms), passes them to a gradient-boosted tree model running on CPU ($<$3ms inference for classical ML), and returns the fraud score. If the score is above a threshold, the transaction is flagged for human review (via a work queue) or auto-declined. The model runs on CPU (not GPU) because gradient boosted trees are highly efficient on CPU, keeping per-request cost low at 20 million trips per day.

\subsection{Content Moderation at Scale (Images / Text)}

Content moderation uses embedding-based classification: a pre-trained vision model embeds each uploaded image into a vector, and a classifier (trained on labeled safe/unsafe examples) scores the vector. Known violating content is stored in a ``perceptual hash'' database (Photo DNA, custom FAISS index of known-bad embeddings); new uploads are matched against this index to detect re-uploads of previously removed content even after minor edits. The serving architecture uses GPU inference for the embedding step (vision transformers are GPU-bound) and CPU for the lightweight classifier. Low-confidence predictions are routed to a human review queue; high-confidence safe/unsafe predictions are acted on automatically.

%--- SECTION 10: CURIOSITY CORNER ---
\section{Curiosity Corner}
\addcontentsline{toc}{section}{Curiosity Corner}

\begin{curiobox}{What is training-serving skew and why is it so hard to detect?}
Training-serving skew is when a model receives different feature values at serving time than it was trained on, causing predictions to be systematically wrong. It is hard to detect because the model does not crash---it produces valid-looking numbers. The degradation is gradual and often correlated with edge cases (users who transact at midnight, new user cohorts, seasonal patterns) rather than uniform across all users. Detection requires monitoring feature distributions at serving time and comparing them to the training distribution (feature drift), monitoring output score distributions (prediction drift), and most importantly, monitoring downstream business metrics (conversion rate, fraud rate) with statistical significance testing against a holdout group. The right fix is a feature store that guarantees the same computation logic in both paths.
\end{curiobox}

\begin{curiobox}{What is RAG and how does it use vector databases?}
Retrieval-Augmented Generation (RAG) is the dominant architecture for grounding LLM responses in private or up-to-date knowledge. The setup: a knowledge base (documents, wiki pages, code, PDFs) is embedded using a text embedding model and stored in a vector database. At query time, the user's question is also embedded; the vector database performs ANN search to find the K most relevant documents (``retrieve''). Those documents are prepended to the LLM's prompt as context (``augment''), and the LLM generates a response conditioned on both the question and the retrieved context (``generate''). This allows LLMs to answer questions about private or recent data without fine-tuning---fine-tuning updates the model's weights and is expensive; RAG updates the knowledge base (a simple vector upsert) and is cheap. The quality of a RAG system is dominated by retrieval quality: if the wrong chunks are retrieved, the LLM will hallucinate regardless of how good the model is.
\end{curiobox}

\begin{curiobox}{What is model drift and how do you detect it in production?}
Model drift is the gradual degradation of a model's predictive accuracy as the real-world data distribution shifts away from the training distribution. There are two subtypes: covariate shift (the input feature distribution changes, e.g., users start using the app differently) and concept drift (the relationship between features and the target changes, e.g., fraud patterns evolve). Detection approaches: (1) Feature drift monitoring: statistical tests (KL divergence, PSI---Population Stability Index) compare the current feature distribution against the training distribution. A PSI $>$ 0.2 signals significant drift and triggers retraining. (2) Prediction drift: monitor the model output score distribution---if the distribution of fraud scores shifts, the model may have drifted. (3) Ground truth monitoring: compare predictions against observed outcomes (did the transaction that was scored ``low fraud risk'' actually result in a chargeback?). Ground truth monitoring is the gold standard but has latency---chargebacks arrive days to weeks after the transaction.
\end{curiobox}

\begin{curiobox}{Why does LLM inference have such different latency characteristics than traditional ML?}
Traditional ML models (gradient boosted trees, logistic regression, even shallow neural nets) operate on structured tabular features: a vector of floats goes in, a scalar prediction comes out. The computation graph is fixed and tiny---microseconds on CPU for gradient boosted trees. LLMs generate tokens \textit{autoregressively}: each token depends on all previous tokens, so generating a 500-token response requires 500 forward passes through a model with billions of parameters. The first token (TTFT) is the most latency-sensitive because the user sees nothing until it arrives. After the first token, throughput (tokens/second) matters more because the user sees the response streaming in. Batching requests together on a GPU dramatically improves throughput (more tokens/second across all users) but increases TTFT for individual requests. vLLM's PagedAttention manages the KV cache (which stores intermediate attention states that grow linearly with context length) as efficiently as an OS manages virtual memory pages, enabling 2--4x throughput gains over naive implementations.
\end{curiobox}

\begin{curiobox}{What is the difference between fine-tuning and RAG, and when do you use each?}
Both fine-tuning and RAG are ways to customize LLM behavior for a specific domain. Fine-tuning updates the model's weights by training on domain-specific examples---the model ``learns'' the domain. It requires labeled examples, compute, and retraining whenever the domain knowledge changes. RAG retrieves relevant context at query time without changing the model---the model ``reads'' the domain. The decision rule is simple: if knowledge changes frequently (news, product inventory, internal wikis), use RAG because updating a vector database is cheap and instant. If the customization is about style, format, or behavior patterns that don't change (``always respond in the tone of our brand'', ``format responses as JSON''), use fine-tuning. Most production systems use both: a fine-tuned model optimized for the task domain, with RAG providing dynamic up-to-date factual grounding.
\end{curiobox}

\begin{curiobox}{How does A/B testing work for ML models in production?}
A/B testing for ML models requires deterministic user assignment: the same user must always see the same model version within an experiment, otherwise their experience is inconsistent and statistical analysis is confounded by the same user being in both groups. The standard implementation uses a hash of the user ID (and experiment ID): \texttt{hash(user\_id + experiment\_id) \% 100 < traffic\_percentage} routes to the treatment model. Experiment results are evaluated on downstream business metrics (click-through rate, conversion, fraud rate) rather than model metrics (AUC, precision), because model metric improvements do not always translate to business metric improvements. Statistical significance testing (typically t-test or Mann-Whitney U for non-normal distributions) determines when the experiment has collected enough samples to make a decision---usually 1--2 weeks for high-traffic surfaces, potentially months for low-traffic ones. Shadow testing (serving the new model's predictions alongside the current model but not using them) is a prerequisite for high-stakes models: it validates that the new model produces sensible predictions in production traffic before any user sees them.
\end{curiobox}

%--- SELF-QUIZ ---
\newpage
\section{Self-Quiz: Interviewer Probes}
\addcontentsline{toc}{section}{Self-Quiz}

\begin{quizbox}{1}
You are designing a music recommendation system for 600 million users. Walk me through the infrastructure components you need beyond a standard database. What is the data flow from a user listening to a song to an updated recommendation?
\end{quizbox}

\begin{quizbox}{2}
What is a feature store and why can't you just query the database at inference time to compute features?
\end{quizbox}

\begin{quizbox}{3}
A user reports that the ``similar items'' feature on your e-commerce site takes 2 seconds to respond. Walk me through how you diagnose and fix it.
\end{quizbox}

\begin{quizbox}{4}
What is an embedding and why do you need a vector database instead of a regular one?
\end{quizbox}

\begin{quizbox}{5}
Your fraud detection model had 98\% accuracy in testing but only 85\% in production after one month. What is most likely happening and how do you fix it?
\end{quizbox}

\newpage
\section*{Model Answers}
\addcontentsline{toc}{section}{Model Answers}

\begin{answerbox}{1}
Beyond a database you need: (1) An event streaming pipeline (Kafka) to capture listening events in real time. (2) A feature pipeline---batch (Spark) for long-horizon features (30-day listening history, genre affinity) writing to an offline store (S3), and streaming (Flink) for short-horizon features (current session) writing to an online store (Redis). (3) A feature store to ensure training and serving use the same computation logic. (4) A vector database storing user embeddings and song embeddings for ANN search. (5) Model serving infrastructure for the ranking model. (6) An A/B testing framework for controlled rollouts. Data flow: user listens to song → event written to Kafka → Flink updates online store features → at next homepage load, serving API fetches features from Redis, runs ANN search against item embedding index, scores top-N candidates with the ranking model → updated recommendations returned $<$100ms.
\end{answerbox}

\begin{answerbox}{2}
A feature store is a centralized system for defining, computing, storing, and serving ML features consistently between training and serving. You can't just query the database at inference time for two reasons: (1) Training-serving skew. The database query you write for serving might not produce the same values as the aggregation used during training---different window definitions, different join logic, missing events in the database that were in the training data lake. The model was trained on one distribution; you're serving it another. (2) Latency. Complex feature computation (rolling window aggregations, user history joins) takes tens to hundreds of milliseconds in a database query. Serving features from a Redis-backed online store takes under 5ms. The feature store's offline store handles training data generation via batch jobs; the online store handles low-latency serving via pre-materialized feature values.
\end{answerbox}

\begin{answerbox}{3}
2-second latency for similar items suggests the ANN search is running on a database not optimized for vector search. Diagnosis: add distributed tracing to isolate which step is slow. Likely culprits: (1) No HNSW index built---falling back to exact nearest-neighbor (linear scan). Fix: build an HNSW index. (2) pgvector at scale beyond its practical ceiling ($>$10M vectors)---index is in disk-based storage, not RAM. Fix: migrate to a dedicated vector database (Qdrant, Weaviate) that keeps the HNSW index in memory. (3) Embeddings computed on-the-fly at query time rather than precomputed and stored. Fix: precompute and store all item embeddings; only embed the query item at request time. After fixing, target $<$50ms for the ANN search step and $<$200ms end-to-end.
\end{answerbox}

\begin{answerbox}{4}
An embedding is a dense vector of floating-point numbers (e.g., 1536 floats) that encodes an item's semantic meaning. Items with similar meaning have embeddings that are geometrically close in the high-dimensional space. You need a vector database because traditional databases use indexes optimized for exact match (B-tree) or range queries (sorted keys)---neither supports efficient nearest-neighbor search in 1536-dimensional space. An exact nearest-neighbor search over 100M vectors is a linear scan: comparing the query vector to 100M stored vectors one by one at 1536 dimensions takes hundreds of milliseconds. A vector database uses HNSW (Hierarchical Navigable Small World graphs) to navigate to approximate nearest neighbors in $O(\log n)$ hops, achieving under 10ms latency at 100M vectors with 99\% recall. The tradeoff: approximate search misses a small fraction of true nearest neighbors, which is acceptable for recommendation and search applications.
\end{answerbox}

\begin{answerbox}{5}
98\% test accuracy dropping to 85\% after one month is a classic training-serving skew or model drift scenario, most likely the latter. First, check for skew: compare the feature distributions at serving time against the training distribution using PSI (Population Stability Index). If features look the same, the model is trained on a stale distribution---the world has changed since the training data was collected (fraud patterns evolved, legitimate transaction patterns shifted). This is concept drift. Fix: retrain the model on the most recent 60--90 days of data. Implement a retraining pipeline that runs weekly or triggered by drift metrics exceeding a threshold. Also add ground-truth monitoring: track the chargeback rate of transactions the model scored as ``low fraud risk''---a rising chargeback rate on low-risk transactions is a direct signal of model degradation. Add feature drift alerts (PSI $>$ 0.2 triggers a Slack alert and initiates automatic retraining).
\end{answerbox}

%--- WHAT'S NEXT ---
\newpage
\section{What's Next}
\addcontentsline{toc}{section}{What's Next}

\begin{insightbox}
\textbf{Layer 1 is complete.} This is the final lesson in the foundational layer. You now have the conceptual bedrock for every system design question you will face: networking (L1-01), push delivery (L1-02), DNS and CDN (L1-03), load balancing (L1-04), relational databases (L1-05), NoSQL (L1-06), caching (L1-07), message queues (L1-08), object storage (L1-09), distributed systems (L1-10), API design (L1-11), security (L1-12), observability (L1-13), and AI/ML infrastructure (L1-14).

Gate 1 is fully unlocked: L1-01 (Networking), L1-02 (APNs/FCM), and L1-08 (Message Queues) are complete. Your first mock interview---the Notification System retry (L3-01)---is now available. The session\_001 9/25 score reflected critical gaps in all three of those areas; they are now closed.

After the Gate 1 mock, Layer 2 Patterns begins. These are cross-cutting design approaches that combine foundation concepts into solutions for classes of problems: fan-out strategies (L2-01), event-driven architecture (L2-02), CQRS (L2-03), the saga pattern (L2-04), sharding strategies (L2-05), replication patterns (L2-06), rate limiting algorithms (L2-07), reliability patterns (L2-08), the notification pipeline pattern (L2-09), and more. Layer 2 is where system design thinking matures from ``I know what a Kafka topic is'' to ``I know when and why to use it over SQS, and what the fan-out topology should look like.''
\end{insightbox}

\end{document}
